% =============================================================================
% related_work_section.tex --- Positional Distillation Paper
% Self-contained related work section (~0.75 page, shortened)
% Usage: % =============================================================================
% related_work_section.tex --- Positional Distillation Paper
% Self-contained related work section (~0.75 page, shortened)
% Usage: % =============================================================================
% related_work_section.tex --- Positional Distillation Paper
% Self-contained related work section (~0.75 page, shortened)
% Usage: % =============================================================================
% related_work_section.tex --- Positional Distillation Paper
% Self-contained related work section (~0.75 page, shortened)
% Usage: \input{related_work_section.tex} in main.tex
% =============================================================================

\section{Related Work}
\label{sec:related_work}

\paragraph{Knowledge Distillation for Language Models.}
Knowledge distillation~\citep{hinton2015distilling,kim2016sequence} for autoregressive LLMs hinges on the choice of divergence and data distribution.
\citet{gu2024minillm} argued for reverse KL, and \citet{agarwal2024onpolicy} introduced on-policy GKD, which supervises student rollouts with teacher logits.
Subsequent work has explored skew KL with adaptive off-policy schedules~\citep{ko2024distillm}, generalized $f$-divergences~\citep{wen2023fdivergence}, mode- vs.\ mean-seeking dynamics~\citep{wu2025rethinking}, and speculative-decoding-style filtering~\citep{xu2025speculative}.
We build on the on-policy reverse-KL paradigm but ask an orthogonal question: holding divergence and data distribution fixed, \emph{which token positions along the student's rollout carry useful signal?}

\paragraph{Token-Level Importance in Distillation and Reasoning.}
Not all tokens contribute equally to learning.
\citet{wang2025beyond8020} showed that $\sim$20\% high-entropy ``forking tokens'' suffice for RL gradients; \citet{vassoyan2025ignorekl} relaxed KL on critical tokens to enable exploration; \citet{li2025preplan} identified a preplan-and-anchor mechanism in attention; and \citet{singh2026functional} pruned reasoning chains by functional importance.
In distillation specifically, recent work selects tokens via teacher verification~\citep{huang2025selectkd}, adaptive temperature~\citep{xie2026adakd}, student entropy along position/class/sample axes~\citep{tavor2026rethinking}, and entropy plus preference ranking~\citep{kim2026tsdkd}.
Our Section~\ref{sec:token_selection} ablation contradicts the shared premise: selecting tokens by top-KL, various top-entropy variants, or soft-weighted joint entropy all underperform \method{} and most underperform plain full-sequence training.
Position is the load-bearing axis; scalar saliencies are not.
This also reconciles \citet{wang2025beyond8020} with our findings: in on-policy distillation, high-entropy tokens are disproportionately concentrated in the uncontaminated early window.

\paragraph{Training Stability and Sequence Length.}
Full-sequence on-policy distillation is unstable: likelihood training degenerates into bland, repetitive text~\citep{holtzman2020curious}, and self-training drives model collapse~\citep{shumailov2024collapse}.
\citet{jang2026veto} proposed a geometric logit-space bridge to suppress harmful gradients on low-confidence tokens.
On the length axis, \citet{chen2025distillingessence} showed that supervised training on the first 50\% of tokens retains $\sim$91\% of full-sequence performance, and \citet{yan2026dasd} addressed exposure bias in long-CoT distillation via temperature scheduling.
Our approach is complementary: restricting the loss scope to early positions, where teacher--student agreement is naturally highest, yields stability through a simpler mechanism and eliminates the degenerate repetition unique to on-policy training.

\paragraph{Curriculum Learning.}
Curriculum learning~\citep{bengio2009curriculum} and scheduled sampling~\citep{bengio2015scheduled} gradually increase task difficulty.
For LLM pretraining, sequence-length warmup stabilizes training~\citep{li2022stability} and variable-length curricula accelerate convergence~\citep{pouransari2024dataset}.
Our progressive position schedule applies the curriculum principle to the distillation loss scope rather than to the input length---a form of curriculum specific to knowledge distillation.

\paragraph{Reasoning Distillation.}
Chain-of-thought prompting~\citep{wei2022chainofthought} made intermediate reasoning the unit of transfer, and DeepSeek-R1~\citep{deepseekr1} showed that SFT-distilling a strong reasoner into a smaller model is often competitive with direct RL.
On reasoning tasks (MATH-500, HumanEval/MBPP, BFCL) and a 1.5B-math-student $\to$ 8B-teacher sweep, \method{} surpasses the teacher on several cells, indicating that gains beyond the teacher are possible when the mode-seeking reverse-KL loss is restricted to the planning region---explaining why the on-policy, early-stopped-rollout regime is more efficient and stable than full-sequence on-policy distillation, which frequently degrades on coding and math.
 in main.tex
% =============================================================================

\section{Related Work}
\label{sec:related_work}

\paragraph{Knowledge Distillation for Language Models.}
Knowledge distillation~\citep{hinton2015distilling,kim2016sequence} for autoregressive LLMs hinges on the choice of divergence and data distribution.
\citet{gu2024minillm} argued for reverse KL, and \citet{agarwal2024onpolicy} introduced on-policy GKD, which supervises student rollouts with teacher logits.
Subsequent work has explored skew KL with adaptive off-policy schedules~\citep{ko2024distillm}, generalized $f$-divergences~\citep{wen2023fdivergence}, mode- vs.\ mean-seeking dynamics~\citep{wu2025rethinking}, and speculative-decoding-style filtering~\citep{xu2025speculative}.
We build on the on-policy reverse-KL paradigm but ask an orthogonal question: holding divergence and data distribution fixed, \emph{which token positions along the student's rollout carry useful signal?}

\paragraph{Token-Level Importance in Distillation and Reasoning.}
Not all tokens contribute equally to learning.
\citet{wang2025beyond8020} showed that $\sim$20\% high-entropy ``forking tokens'' suffice for RL gradients; \citet{vassoyan2025ignorekl} relaxed KL on critical tokens to enable exploration; \citet{li2025preplan} identified a preplan-and-anchor mechanism in attention; and \citet{singh2026functional} pruned reasoning chains by functional importance.
In distillation specifically, recent work selects tokens via teacher verification~\citep{huang2025selectkd}, adaptive temperature~\citep{xie2026adakd}, student entropy along position/class/sample axes~\citep{tavor2026rethinking}, and entropy plus preference ranking~\citep{kim2026tsdkd}.
Our Section~\ref{sec:token_selection} ablation contradicts the shared premise: selecting tokens by top-KL, various top-entropy variants, or soft-weighted joint entropy all underperform \method{} and most underperform plain full-sequence training.
Position is the load-bearing axis; scalar saliencies are not.
This also reconciles \citet{wang2025beyond8020} with our findings: in on-policy distillation, high-entropy tokens are disproportionately concentrated in the uncontaminated early window.

\paragraph{Training Stability and Sequence Length.}
Full-sequence on-policy distillation is unstable: likelihood training degenerates into bland, repetitive text~\citep{holtzman2020curious}, and self-training drives model collapse~\citep{shumailov2024collapse}.
\citet{jang2026veto} proposed a geometric logit-space bridge to suppress harmful gradients on low-confidence tokens.
On the length axis, \citet{chen2025distillingessence} showed that supervised training on the first 50\% of tokens retains $\sim$91\% of full-sequence performance, and \citet{yan2026dasd} addressed exposure bias in long-CoT distillation via temperature scheduling.
Our approach is complementary: restricting the loss scope to early positions, where teacher--student agreement is naturally highest, yields stability through a simpler mechanism and eliminates the degenerate repetition unique to on-policy training.

\paragraph{Curriculum Learning.}
Curriculum learning~\citep{bengio2009curriculum} and scheduled sampling~\citep{bengio2015scheduled} gradually increase task difficulty.
For LLM pretraining, sequence-length warmup stabilizes training~\citep{li2022stability} and variable-length curricula accelerate convergence~\citep{pouransari2024dataset}.
Our progressive position schedule applies the curriculum principle to the distillation loss scope rather than to the input length---a form of curriculum specific to knowledge distillation.

\paragraph{Reasoning Distillation.}
Chain-of-thought prompting~\citep{wei2022chainofthought} made intermediate reasoning the unit of transfer, and DeepSeek-R1~\citep{deepseekr1} showed that SFT-distilling a strong reasoner into a smaller model is often competitive with direct RL.
On reasoning tasks (MATH-500, HumanEval/MBPP, BFCL) and a 1.5B-math-student $\to$ 8B-teacher sweep, \method{} surpasses the teacher on several cells, indicating that gains beyond the teacher are possible when the mode-seeking reverse-KL loss is restricted to the planning region---explaining why the on-policy, early-stopped-rollout regime is more efficient and stable than full-sequence on-policy distillation, which frequently degrades on coding and math.
 in main.tex
% =============================================================================

\section{Related Work}
\label{sec:related_work}

\paragraph{Knowledge Distillation for Language Models.}
Knowledge distillation~\citep{hinton2015distilling,kim2016sequence} for autoregressive LLMs hinges on the choice of divergence and data distribution.
\citet{gu2024minillm} argued for reverse KL, and \citet{agarwal2024onpolicy} introduced on-policy GKD, which supervises student rollouts with teacher logits.
Subsequent work has explored skew KL with adaptive off-policy schedules~\citep{ko2024distillm}, generalized $f$-divergences~\citep{wen2023fdivergence}, mode- vs.\ mean-seeking dynamics~\citep{wu2025rethinking}, and speculative-decoding-style filtering~\citep{xu2025speculative}.
We build on the on-policy reverse-KL paradigm but ask an orthogonal question: holding divergence and data distribution fixed, \emph{which token positions along the student's rollout carry useful signal?}

\paragraph{Token-Level Importance in Distillation and Reasoning.}
Not all tokens contribute equally to learning.
\citet{wang2025beyond8020} showed that $\sim$20\% high-entropy ``forking tokens'' suffice for RL gradients; \citet{vassoyan2025ignorekl} relaxed KL on critical tokens to enable exploration; \citet{li2025preplan} identified a preplan-and-anchor mechanism in attention; and \citet{singh2026functional} pruned reasoning chains by functional importance.
In distillation specifically, recent work selects tokens via teacher verification~\citep{huang2025selectkd}, adaptive temperature~\citep{xie2026adakd}, student entropy along position/class/sample axes~\citep{tavor2026rethinking}, and entropy plus preference ranking~\citep{kim2026tsdkd}.
Our Section~\ref{sec:token_selection} ablation contradicts the shared premise: selecting tokens by top-KL, various top-entropy variants, or soft-weighted joint entropy all underperform \method{} and most underperform plain full-sequence training.
Position is the load-bearing axis; scalar saliencies are not.
This also reconciles \citet{wang2025beyond8020} with our findings: in on-policy distillation, high-entropy tokens are disproportionately concentrated in the uncontaminated early window.

\paragraph{Training Stability and Sequence Length.}
Full-sequence on-policy distillation is unstable: likelihood training degenerates into bland, repetitive text~\citep{holtzman2020curious}, and self-training drives model collapse~\citep{shumailov2024collapse}.
\citet{jang2026veto} proposed a geometric logit-space bridge to suppress harmful gradients on low-confidence tokens.
On the length axis, \citet{chen2025distillingessence} showed that supervised training on the first 50\% of tokens retains $\sim$91\% of full-sequence performance, and \citet{yan2026dasd} addressed exposure bias in long-CoT distillation via temperature scheduling.
Our approach is complementary: restricting the loss scope to early positions, where teacher--student agreement is naturally highest, yields stability through a simpler mechanism and eliminates the degenerate repetition unique to on-policy training.

\paragraph{Curriculum Learning.}
Curriculum learning~\citep{bengio2009curriculum} and scheduled sampling~\citep{bengio2015scheduled} gradually increase task difficulty.
For LLM pretraining, sequence-length warmup stabilizes training~\citep{li2022stability} and variable-length curricula accelerate convergence~\citep{pouransari2024dataset}.
Our progressive position schedule applies the curriculum principle to the distillation loss scope rather than to the input length---a form of curriculum specific to knowledge distillation.

\paragraph{Reasoning Distillation.}
Chain-of-thought prompting~\citep{wei2022chainofthought} made intermediate reasoning the unit of transfer, and DeepSeek-R1~\citep{deepseekr1} showed that SFT-distilling a strong reasoner into a smaller model is often competitive with direct RL.
On reasoning tasks (MATH-500, HumanEval/MBPP, BFCL) and a 1.5B-math-student $\to$ 8B-teacher sweep, \method{} surpasses the teacher on several cells, indicating that gains beyond the teacher are possible when the mode-seeking reverse-KL loss is restricted to the planning region---explaining why the on-policy, early-stopped-rollout regime is more efficient and stable than full-sequence on-policy distillation, which frequently degrades on coding and math.
 in main.tex
% =============================================================================

\section{Related Work}
\label{sec:related_work}

\paragraph{Knowledge Distillation for Language Models.}
Knowledge distillation~\citep{hinton2015distilling,kim2016sequence} for autoregressive LLMs hinges on the choice of divergence and data distribution.
\citet{gu2024minillm} argued for reverse KL, and \citet{agarwal2024onpolicy} introduced on-policy GKD, which supervises student rollouts with teacher logits.
Subsequent work has explored skew KL with adaptive off-policy schedules~\citep{ko2024distillm}, generalized $f$-divergences~\citep{wen2023fdivergence}, mode- vs.\ mean-seeking dynamics~\citep{wu2025rethinking}, and speculative-decoding-style filtering~\citep{xu2025speculative}.
We build on the on-policy reverse-KL paradigm but ask an orthogonal question: holding divergence and data distribution fixed, \emph{which token positions along the student's rollout carry useful signal?}

\paragraph{Token-Level Importance in Distillation and Reasoning.}
Not all tokens contribute equally to learning.
\citet{wang2025beyond8020} showed that $\sim$20\% high-entropy ``forking tokens'' suffice for RL gradients; \citet{vassoyan2025ignorekl} relaxed KL on critical tokens to enable exploration; \citet{li2025preplan} identified a preplan-and-anchor mechanism in attention; and \citet{singh2026functional} pruned reasoning chains by functional importance.
In distillation specifically, recent work selects tokens via teacher verification~\citep{huang2025selectkd}, adaptive temperature~\citep{xie2026adakd}, student entropy along position/class/sample axes~\citep{tavor2026rethinking}, and entropy plus preference ranking~\citep{kim2026tsdkd}.
Our Section~\ref{sec:token_selection} ablation contradicts the shared premise: selecting tokens by top-KL, various top-entropy variants, or soft-weighted joint entropy all underperform \method{} and most underperform plain full-sequence training.
Position is the load-bearing axis; scalar saliencies are not.
This also reconciles \citet{wang2025beyond8020} with our findings: in on-policy distillation, high-entropy tokens are disproportionately concentrated in the uncontaminated early window.

\paragraph{Training Stability and Sequence Length.}
Full-sequence on-policy distillation is unstable: likelihood training degenerates into bland, repetitive text~\citep{holtzman2020curious}, and self-training drives model collapse~\citep{shumailov2024collapse}.
\citet{jang2026veto} proposed a geometric logit-space bridge to suppress harmful gradients on low-confidence tokens.
On the length axis, \citet{chen2025distillingessence} showed that supervised training on the first 50\% of tokens retains $\sim$91\% of full-sequence performance, and \citet{yan2026dasd} addressed exposure bias in long-CoT distillation via temperature scheduling.
Our approach is complementary: restricting the loss scope to early positions, where teacher--student agreement is naturally highest, yields stability through a simpler mechanism and eliminates the degenerate repetition unique to on-policy training.

\paragraph{Curriculum Learning.}
Curriculum learning~\citep{bengio2009curriculum} and scheduled sampling~\citep{bengio2015scheduled} gradually increase task difficulty.
For LLM pretraining, sequence-length warmup stabilizes training~\citep{li2022stability} and variable-length curricula accelerate convergence~\citep{pouransari2024dataset}.
Our progressive position schedule applies the curriculum principle to the distillation loss scope rather than to the input length---a form of curriculum specific to knowledge distillation.

\paragraph{Reasoning Distillation.}
Chain-of-thought prompting~\citep{wei2022chainofthought} made intermediate reasoning the unit of transfer, and DeepSeek-R1~\citep{deepseekr1} showed that SFT-distilling a strong reasoner into a smaller model is often competitive with direct RL.
On reasoning tasks (MATH-500, HumanEval/MBPP, BFCL) and a 1.5B-math-student $\to$ 8B-teacher sweep, \method{} surpasses the teacher on several cells, indicating that gains beyond the teacher are possible when the mode-seeking reverse-KL loss is restricted to the planning region---explaining why the on-policy, early-stopped-rollout regime is more efficient and stable than full-sequence on-policy distillation, which frequently degrades on coding and math.
